\documentclass[a4paper,10pt]{article}
\usepackage[utf8]{inputenc}
\usepackage[T1]{fontenc}
\usepackage[french]{babel}
\usepackage{amsmath,amssymb}
\usepackage{geometry}
\usepackage{array,tabularx}
\usepackage{tikz}
\geometry{margin=1.6cm}
\pagestyle{empty}
\newcommand{\ligne}{\par\vspace{0.55cm}\noindent\dotfill\par}
\newcommand{\carreaux}[1]{\begin{center}\begin{tikzpicture}\draw[step=0.5cm,gray!35,very thin] (0,0) grid (17,#1);\end{tikzpicture}\end{center}}
\newenvironment{exercice}[2]{\paragraph{Exercice #1 \hfill #2 points}\mbox{}\par}{\medskip}
\begin{document}
\begin{center}\Large\bfseries Devoir surveillé 15 - Agrandissement et réduction\end{center}
\vspace{2mm}
\begin{exercice}{1}{4}Une figure est agrandie avec un coefficient $1,5$. Une longueur de $8$ cm devient-elle ?\end{exercice}\begin{exercice}{2}{5}Un triangle de côtés $4$, $6$, $7$ cm est réduit avec un coefficient $0,75$. Donner les nouvelles longueurs.\end{exercice}\begin{exercice}{3}{5}Une aire est multipliée par $9$. Quel peut être le coefficient d'agrandissement des longueurs ?\end{exercice}\begin{exercice}{4}{6}Un solide est agrandi avec un coefficient $2$. Comment évoluent ses longueurs, ses aires et son volume ?\end{exercice}\end{document}
